\documentclass[12pt,a4paper]{article}

\usepackage{fontawesome5}
\providecommand{\faPencilAlt}{\faIcon{pencil-alt}}

\usepackage{styles/couleurs}
\usepackage{styles/mise_en_page}
\usepackage{styles/boites}
\usepackage{styles/maths}
\usepackage{styles/figures}
\usepackage{styles/evaluations}

\usepackage{enumitem}
\usepackage{array}
\usepackage{tabularx}
\usepackage{amsmath}
\usepackage{tikz}

\setlist[itemize]{leftmargin=1.4em,itemsep=0.15em,topsep=0.2em}
\setlist[enumerate]{leftmargin=*,itemsep=0.25em,topsep=0.25em}

\begin{document}

\enteteinterro[Calculatrice autorisée]{4e}{Statistiques et diagrammes}{15 min}{12}

\exo{4}[Calculer]
Une enquête porte sur \(40\) élèves.

\[
\begin{array}{c|cccc}
\text{Choix} & A & B & C & D\\
\hline
\text{Effectif} & 16 & 10 & 8 & 6
\end{array}
\]

Calcule le pourcentage correspondant à chaque choix. \pts{4}

\lignes{5}

\exo{4}[Calculer]
Calcule l'angle du secteur correspondant à chaque choix. \pts{4}

\lignes{5}

\exo{2}[Raisonner]
Un diagramme annonce \(42\%\), \(31\%\), \(18\%\) et \(14\%\).
Peut-il représenter une répartition complète ? Justifie. \pts{2}

\lignes{3}

\exo{2}[Communiquer]
Cite deux informations qui doivent accompagner un diagramme
pour qu'il soit correctement interprété. \pts{2}

\lignes{2}

\end{document}
